\section{Evaluation}
\subsection{Experiments and Results}
% We target the AMD Alveo U250 FPGA (XCU250) for fair comparison with prior work. 
We target the AMD Alveo U280 FPGA (XCU280), part xcu280-fsvh2892-2L-e. 
We build large multipliers for sizes ranging from $160$ to $1024$ bits using exact tiling with $77$-bit base or non-exact tiling with $43$-bit base multipliers as discussed in Section~\ref{sec_tiler} and used $64$-bit AMD IP blocks for $1024$-bit multiplier.
The exact tiling is the same as the non-exact tiling for the $1024$-bit multiplier using the $64$-bit AMD IP (four levels of $2\times2$ grids).
Additionally, the $77$-bit multipliers have fewer combinations as the exact tiling result in two remaining rectangular areas that are not Karatsuba friendly.
Wide additions are pipelined with segments of size $64$ bits using the lowest-latency mix of the lockstep and interleaved strategies described in Section~\ref{sec_fast_pipelined_comb_logic}.
Our analysis compares these implementations of the $1024$-bit multiplier against previous work in terms of frequency, LUT utilization, and DSP utilization.
Additionally, a unified metric that is used by IMpress~\cite{zhangIMpressLargeInteger2022} called Normalized Area will be used.
Alveo U280 has $1082$K LUTs and $8490$ DSPs~\cite{U280SwitchingCharacteristics}, giving a ratio of around $126$ LUTs per DSP. 
The normalized area is defined as:
\begin{equation*}
\text{Normalized Area} = \#\text{LUTs} + (126 \times \#\text{DSPs})
\end{equation*}
where each DSP is considered equivalent to 126 LUTs.

Table ~\ref{table_decomposition_schemes} provides a detailed comparison of different decomposition schemes for various target and base multiplier sizes. 
Figure~\ref{fig_bar_comparison_dsp_area} depicts the DSP and normalized area utilization for different decomposition schemes using our $77$-bit and $43$-bit base multipliers.
Consistently across all target sizes, using more Karatsuba levels results in lower DSP utilization and overall normalized area while achieving higher frequencies most of the time.
The $43$-bit base implementation has better hierarchical nature, and thus it is more superior when using more Karatsuba levels while the $77$-bit base multiplier generally outperforms it when using fewer Karatsuba levels.

% comparing to the other work
Table ~\ref{table_decomposition_schemes} highlights the best configuration per implementation in bold.
Our IP-based multipliers consistently achieved significantly higher frequencies than other work but being less efficient in terms of resource utilization due to the inefficient $64$-bit IP blocks. 
Our $43$-bit achieved the same high frequency while utilizing fewer DSPs, and using comparable LUTs.
% Figure explanation
Figure~\ref{fig_dualaxis_best_diff1024} compares the models with the least number of DSPs across implementations in terms of the maximum achieved frequency against their normalized area and DSP usage.
The red points denote the DSP usage axis while the blue points denote the normalized area axis. 
The HLS implementation achieves the lowest DSP usage but at the cost of low frequency and high normalized area.
IMpress1 and IMpress2 achieve moderate DSP and normalized area utilization but also suffer from low frequency.
Our IP-based implementation and  $77$-bit primitive-based implementation achieve a  higher frequency than other implementations but at the cost of higher DSP and normalized area.
Our 43-bit primitive-based implementation (the circle) achieves the highest frequency while using moderate DSP and normalized area utilization compared to other implementations.
It achieved around 83\% higher frequency than IMpress1 implementation while using 16\% fewer LUTs and the same number of DSPs.
Additionally, it achieved around 56\% higher frequency than IMpress2 while using 17\% fewer DSPs and slightly less number of LUTs. 
This shows that coupling the tiling automation with an efficient hand-crafted base primitive is key to maximizing throughput and resource efficiency.

%%%%% Detailed comparison of different decomposition schemes for various multiplier sizes %%%%

\begin{table}[t]
\centering
\caption{Detailed comparison of different decomposition schemes for various multiplier sizes. Exact tiling is used with the $77\times77$ base multiplier, while non-exact tiling is used with the $43\times43$ base multiplier. Bold indicates the best configuration per implementation.}
\resizebox{1\linewidth}{!}{
\begin{tabular}{cccccccc}
\toprule
\textbf{Target Size} & \textbf{Base} & 
\shortstack{\textbf{\#Karat.} \\ [-0.5ex] \textbf{Levels}} & \textbf{LUTs} & \textbf{FFs} & \textbf{DSPs} & 
\shortstack{\textbf{Max Freq.} \\ [-0.5ex] \textbf{(MHz)}} & 
\shortstack{\textbf{Latency} \\ [-0.5ex] \textbf{(Cycles)}} \\
\midrule
160 & primitive & 0 & 3461 & 6112 & 61 & 830 & 35 \\
    & 77 & 1 & 3770 & 7507 & 62 & 855 & 31 \\
    & 43 & 2 & 5141 & 9202 & 45 & 862 & 30 \\
\midrule
256 & primitive & 0 & 9947 & 15710 & 161 & 760 & 64 \\
     & 77 & 1 & 8473 & 17718 & 114 & 776 & 40 \\
     & 43 & 2 & 10858 & 20171 & 90 & 842 & 35 \\
\midrule
384 & 77 & 2 & 21350 & 38330 & 279 & 583 & 56 \\
     & 43 & 2 & 23208 & 43803 & 180 & 751 & 48 \\
\midrule
512 &  77 & 2 & 33122 & 61525 & 370 & 619 & 72 \\
     & 43 & 3 & 38131 & 68870 & 270 & 629 & 60 \\
\midrule
1024 & 77 & 3 & 123472 & 207765 & 1220 & 478 & 123 \\
     & 43 & 4 & 128252 & 221612 & 810 & 458 & 100 \\
\midrule
1024 & IP & 4 & 105717 & 169858 & 1296 & 459 & 89 \\
\midrule 
1024 & HLS & 6 & 165618 & N/A & 729 & 298 & 24 \\
     & OpenNTT \cite{sinharoyOpenNTTAutomatedToolchain2025} & N/A & 12717 & 14217 & 80 & 320 & 689 \\
    & IMpress1 \cite{zhangIMpressLargeInteger2022} & N/A & 150400 & N/A & 810 & 250 & 23 \\
     & IMpress2 \cite{zhangIMpressLargeInteger2022} & N/A & 129338 & N/A & 945 & 295 & 22 \\
     & BeamSearch \cite{kummTilingHeuristics2020} & N/A & 133975 & N/A & 1471 & 283 & N/A \\
\bottomrule
\end{tabular}
}
\label{table_decomposition_schemes}
\end{table}



%%% figures for different implementations comparison %%%
%Utilization in different levels of the decomposition
\begin{figure}[t]
     \centering
     \begin{subfigure}{0.5\textwidth}
         \centering
         \includegraphics[width=\textwidth]{figures/evaluation/bar_dsp_usage.pdf}
         \caption{DSP Utilization with Different Decomposition Schemes}
         \label{fig_bar_dsp_usage}
     \end{subfigure}   
     \hfill
     \begin{subfigure}{0.5\textwidth}
         \centering
         \includegraphics[width=\textwidth]{figures/evaluation/bar_normalized_area.pdf}
         \caption{Normalized Area with Different Decomposition Schemes}
         \label{fig_bar_normalized_area}
     \end{subfigure}
     \hfill
     \caption{Resource Utilization for using our 77-bit base (\textsuperscript{*}) and our 43-bit base (\textsuperscript{**}) multipliers, with varying numbers of Karatsuba levels.}
    
      \label{fig_bar_comparison_dsp_area}
   
\end{figure}

%%% end figure for different implementations comparison %%%







\begin{figure}[t]
\centering
\includegraphics[width=\columnwidth]{figures/evaluation/best_scatter_dual_axis_frequency.pdf}
\caption{Area and Performance comparison of different implementations.}
\label{fig_dualaxis_best_diff1024}
\end{figure}

%%%

\subsection{Latency Mitigation}
As shown, our primitive-based large multipliers achieve significantly higher frequencies than prior work, but it suffers from higher latency.
To mitigate this, we divide the wide addition to larger segment sizes (i.e. 128 bits instead of 64 bits).
Table~\ref{table_combiner_size} shows the impact of varying the combiner segment size on the latency and performance of a $1032$-bit multiplier with a $43$-bit base multiplier.
Increasing the segment size reduces the latency significantly, with a trade-off of slightly increased LUT usage and reduced maximum frequency.
Using a 256-bit segment size reduces the latency from 100 to 58 cycles while still achieving higher frequency than other implementations.


%%% table for latency mitigation %%% 

\begin{table}[t]
\centering
\caption{Impact of combiner base size on latency and performance for 1032-bit multiplier 4-Karatsuba with 43-bit base}
\resizebox{0.8\linewidth}{!}{
\begin{tabular}{ccccc}
\toprule
\shortstack{\textbf{Combiner} \\ [-0.5ex] \textbf{Segment Size}} & \textbf{LUTs} & \textbf{FFs} & 
\shortstack{\textbf{Max Freq.} \\ [-0.5ex] \textbf{(MHz)}} & 
\shortstack{\textbf{Latency} \\ [-0.5ex] \textbf{(Cycles)}} \\
\midrule
96 & 127225 & 225670 & 379 & 88 \\
128 & 125480 & 219590 & 359 & 70 \\
192 & 133741 & 216455 & 368 & 66 \\
256 & 136056 & 213833 & 348 & 58 \\
\bottomrule
\end{tabular}
}
\label{table_combiner_size}
\end{table}


